% Generated by task tex-groundtruth from dossiers/gt-flock-ring.md (section 8, findings),
% with the corrections of dossiers/verify-gt-flock-ring.md.
\section{The BLAKE2s validity phase (Flock) and ring switching: findings}\label{sec:gen-flock-ring-findings}

\begin{sloppypar}
Short forms in this section; the number after the colon is the line. leanVM is read at its pin
\code{a386121f}, leanerVM at \code{main} \code{b435631}.
\begin{itemize}
\item leanerVM: \code{bp:} is the blueprint \file{docs/roadmap/protocol-blueprint.md},
  \code{st:} the status \file{docs/roadmap/protocol-status.md}; Lean files are under
  \file{LeanerVM/Protocol/}.
\item Specification, under \file{doc/leanvm/body/}: \code{c-flock:} is
  \file{c-flock-protocol.tex}, \code{a-ring:} \file{a-ring-switching.tex}, \code{03:}
  \file{03-proving-primitives.tex}, \code{06:} \file{06-bus-interactions.tex}, \code{07:}
  \file{07-instruction-tables.tex}, \code{08:} \file{08-end-to-end-protocol.tex}.
\item Rust, under \file{crates/flock/src/}: \code{hash.rs}, \code{gf2.rs}, \code{zerocheck.rs},
  \code{lincheck.rs}; under \file{crates/flock/src/zerocheck/}:
  \code{univariate_skip_optimized.rs}; under \file{crates/lean_vm/src/}: \code{cpu/mod.rs},
  \code{cpu/layout.rs}, \code{hash_flock.rs}, \code{tables.rs}; under \file{crates/pcs/src/}:
  \code{stack_open.rs}, \code{ring_switch.rs}, \code{whir_config.rs}; under
  \file{crates/primitives/src/field/}: \code{gf2_8.rs}, \code{phi8_tower.rs},
  \code{gf2_64x3.rs}; \code{primitives/src/multilinear.rs} is
  \file{crates/primitives/src/multilinear.rs}; \code{transcript.rs} is
  \file{crates/fiat_shamir/src/transcript.rs}.
\item Python: \code{verifier.py} is \file{python-verifier/verifier.py}.
\end{itemize}
\end{sloppypar}

No finding is critical: no Lean theorem about Flock exists yet that could be wrong or vacuous.
The four major findings are all in text that the Flock roadmap will build from.

% ---------------------------------------------------------------- 8.1
\begin{finding}{major}{f:flock-ring-interface-limb-claims}{The Flock interface of Layer 9 consumes
the limb claims and emits one claim; leanVM's Flock reads no claim and the limb claims are opened}
\fhead{Evidence} \Cref{tab:gen-flock-ring-limbs}, and the dossier's section 5.1. The deployed
Flock verifier reads no claim of the pool: the eighteen limb claims stay in the pool and are
opened as evaluation claims on $\qflock$ (strided). Layer 9's \lean{FlockInterface}, with
\lean{StmtIn := ColumnClaims} (the eighteen claims) and \lean{StmtOut := WeightedClaim}, drops
them. Its knowledge soundness is not provable: take $q$ with a valid $\qflock$ and limb claims
with wrong values; \lean{relIn} fails, the deployed verifier never reads the claims, the honest
transcript is accepted and \lean{relOut} holds: error 1. A reduction of Layer 9's shape could be
made sound only by folding the eighteen claims into the weighted claim with a fresh challenge,
which leanVM does not do.

\fhead{Classification} An error of the blueprint, to be fixed; the spine is faithful, Layer 9 and
the tracker's section on the Flock phase (P6) are not.

\fhead{Proposed change}
\begin{change}{Layer 9: the Flock phase at the spine's seams}
\asis (\code{bp:1075-1092}) The structure
\begin{leancode}
structure FlockInterface (I : M3Instance) where
  reduction : OracleReduction []ₒ (StmtIn := ColumnClaims) (OStmtIn := fun _ : Unit ↦ Column μ) Unit
      (StmtOut := WeightedClaim) (OStmtOut := fun _ : Unit ↦ Column μ) Unit pSpecFlock
  limbColumns : Column μ → Fin 18 → Column (s.τ 5)      -- the limbs: strided slots of q_flock
  relIn : Set _   -- the eighteen column claims are true of `limbColumns`, and aux (Flock's R1CS)
  relOut : Set _  -- the weighted claim holds for q
\end{leancode}
\src{docs/roadmap/protocol-blueprint.md:1077-1082 at b435631}
and \enquote{\lean{relIn} carries the strong auxiliary predicate of Layer 3 \dots\
\lean{limbColumns} is the strided reader of Layer 3}. (The citation check noted that the dossier
cites this block as \code{bp:1075-1086}; it reproduces lines 1077 to 1082.)

\asproposed A sketch in the blueprint's style, not compiled:
\begin{leancode}
/-- What the Flock roadmap supplies: the phase at the spine's seams, for an instance that says
where the packed witness is (`FlockRegion`, Layer 3). -/
structure FlockPhase (I : M3Instance) (R : FlockRegion I) where
  phase : Phase.Def I (I.Stmt × PubOut I) (I.Stmt × FlockOut I)
  security : Phase.Security I phase (Seam.pub I) (Seam.flock I)   -- completeness included
  err_le : ∑ i, phase.err i ≤ (4 * R.kBatch + 163) / |E| + 2 ^ 32 / |E|
\end{leancode}
with the text: \enquote{The verifier reads no claim of the pool. It hands the column claims on
unchanged, the eighteen limb claims among them (\code{cpu/mod.rs:790-814}: they are opened as
evaluation claims on $\qflock$), and adds one weighted claim, the ring-switched one. The schedule
is: $k_{\mathrm{batch}} + 1$ challenges; 64 values; one challenge; $8 + k_{\mathrm{batch}}$ rounds
of two values and one challenge; two values; one challenge; 8 rounds of two values and one
challenge; 64 values; six challenges.}

\reason Soundness of the sketch as written fails with error 1.
\end{change}
\end{finding}

% ---------------------------------------------------------------- 8.2
\begin{finding}{major}{f:flock-ring-slot-map-circular}{The slot map of the limb columns is placed
inside the Flock interface, so the instance and the interface each need the other}
\fhead{Evidence} Cannot be stated as written. \code{bp:846-848} \enquote{\code{leanIsaInstance.layout}
combines the two, the slot map being \#3's \code{FlockInterface.limbColumns} (Layer 9)}, and
\code{bp:1077} \lean{structure FlockInterface (I : M3Instance)}. The instance's \lean{layout}
field needs a value of a structure that takes the instance as its parameter. The slot map is
eighteen numbers (\code{hash_flock.rs:96-115}) and a generic identity on multilinear tables.

\fhead{Classification} An error of the blueprint.

\fhead{Proposed change}
\begin{enumerate}[label=(\alph*)]
\item \lean{blake2sLimbSlot : Fin 18 → Fin 256} in \file{LeanerVM/Parameters/}, transcribed from
  its source (the blueprint's Category B: \code{hash_flock.rs:87-115}, cross-checked with
  \code{verifier.py:847-850});
\item the low-index selection identity next to \lean{evalMle_append_boolVec} in
  \file{ToCompPoly/Multilinear.lean}, generic (a table of $m + k$ variables evaluated at
  \lean{boolVec j ++ z} is its low-index slice evaluated at $z$);
\item Layer 3 builds the strided reader from the two;
\item \lean{limbColumns} is deleted from Layer 9 and from the tracker's sections on the adaptor
  (I2) and on the Flock phase (P6).
\end{enumerate}
The boundary paragraph with leanISA (\code{bp:1439-1441}) becomes \enquote{reconciled by the
instance's layout}.
\end{finding}

% ---------------------------------------------------------------- 8.3
\begin{finding}{major}{f:flock-ring-opaque-aux}{The Flock phase cannot be written over an abstract
instance: the instance gives it only an opaque predicate}
\fhead{Evidence} Cannot be stated as written.
\begin{leancode}
  /-- What the Flock phase proves of the committed region it owns, beyond the polynomial
  checks: the statement its honest prover can convince the verifier of, so it names the
  committed witness (Flock's R1CS on `q_flock`), never only a consequence of it. -/
  aux : Column μ → Prop
  /-- The auxiliary predicate is decidable, so that the relation is. -/
  decAux : DecidablePred aux
\end{leancode}
\src{LeanerVM/Protocol/Spine/Instance.lean:143-148 at b435631}
Layer 9 is \enquote{module, over \lean{I}} (\code{bp:1072}), and the convention \emph{The wall}
(\code{bp:331}) lists the modules allowed to know leanISA; the Flock phase is not among them. But
a phase that knows only \lean{I.aux : Column μ → Prop} cannot name the column, the number of
compressions or the relation it is to prove; the sketch itself reaches for \lean{s.τ 5}, a
leanISA size. For an abstract \lean{I} the only phases one can write are those of the probe (the
dossier's section 9.1: a Flock phase that checks nothing is knowledge sound when \lean{aux} is
trivial).

\fhead{Classification} An error of the blueprint (a design gap the spine left open).

\fhead{Proposed change} A sketch: Layer 3 supplies, and Layer 9 takes,
\begin{leancode}
/-- Where the packed witness of the BLAKE2s circuit sits, and what `aux` says of it. -/
structure FlockRegion (I : M3Instance) where
  col : I.ColumnId                        -- the column q_flock
  kBatch : ℕ                              -- log₂ of the number of compressions
  height : I.κ col = 8 + kBatch           -- 2^8 packed words per compression, low coordinates
  aux_iff : ∀ q, I.aux q ↔ Flock.Holds kBatch (I.column q col)   -- R1CS and constant position
\end{leancode}
\lean{Flock.Holds} is the Flock roadmap's predicate on a packed column and mentions no instance.
The toy keeps \lean{aux := True} and has no \lean{FlockRegion}. The wall stands: the phase imports
nothing from the arithmetization.
\end{finding}

% ---------------------------------------------------------------- 8.4
\begin{finding}{major}{f:flock-ring-test20-completeness}{Acceptance test 20 files a completeness
failure under a soundness error, and attributes to the protocol an inverse that only the Rust
prover takes}
\fhead{Evidence} The dossier's section 4. No verifier takes the inverse of a value that depends on
a challenge: both derive the untransmitted coefficient of a round polynomial by a multiplication
(\code{transcript.rs:297-302}; \code{verifier.py:411-413}). (The citation check corrected the
dossier's sentence \enquote{The verifiers take no inverse}, which is literally false: both
verifiers invert constants, for example \code{univariate_skip_optimized.rs:96-99},
\code{primitives/src/multilinear.rs:185}, \code{verifier.py:1099} and \code{:1110}.) The Rust
prover takes one, to avoid computing $G(0)$:
\begin{srccode}
fn send_round(ps: &mut impl Transmitter, claim: F192, r_eq: F192, g1: F192, g_inf: F192, chis: &mut Vec<F192>) -> F192 {
    let g0 = (claim + r_eq * g1) * (F192::ONE + r_eq).inv();
    ps.add_round_poly(&[g0, g0 + g1 + g_inf, g_inf], true);
    let chi = ps.sample();
    chis.push(chi);
    // G(X) = G(0)·(1+X) + G(1)·X + G(inf)·X·(1+X).
    g0 + chi * (g0 + g1 + (F192::ONE + chi) * g_inf)
}
\end{srccode}
\src{crates/flock/src/zerocheck.rs:116-123}
and the inverse of zero is zero, by contract (\code{gf2_64x3.rs:137}). At $r_{\mathrm{eq}} = 1$
the true claim is $(1 + 1)·G(0) + 1·G(1) = G(1)$: it carries no information on $G(0)$. The prover
computes $g_0 = 0$, sends $c_1 = G(1) + G(∞)$ (wrong unless $G(0) = 0$) and the right $c_2$; the
verifier derives $c_0 = 0$, and the error propagates into the derived $v_c$. The verifier rejects,
at the lincheck terminal identity; no panic on either side. The challenges concerned are the
$k_{\mathrm{batch}} + 1$ sampled equality coordinates only, with probability at most
$(k_{\mathrm{batch}} + 1)/|\E| ≤ 33/2^{192}$, times the chance that $G(0) ≠ 0$. It is not the
protocol's: a prover that sends the true $c_1, c_2$ is accepted at $r_{\mathrm{eq}} = 1$, and the
specification never mentions an inverse (\code{c-flock:69}, \code{03:110}). Derived by reading and
confirmed on a model of the two formulas over the pinned Python field (the dossier's probe 9.2);
not executed in Rust. The spine has no place for a completeness error at all
(\lean{Component.Security} extends \lean{Complete}, which is perfect completeness).

\fhead{Classification} An error of the blueprint.

\fhead{Proposed change}
\begin{change}{acceptance test 20}
\asis (\code{bp:1323-1325}) \enquote{20. \textbf{Perfect completeness needs no
exceptional-challenge clause} in Layers 4--8 and 10: no inverse of a challenge is taken. Flock's
$(1 + r_{\mathrm{eq}})^{-1}$ at $r_{\mathrm{eq}} = 1$ is \#3's, inside \lean{flockError}. Witness:
\lean{piop_perfectCompleteness} has no side condition on challenges.}

\asproposed \enquote{20. \textbf{Perfect completeness needs no exceptional-challenge clause, in
any phase.} No verifier takes the inverse of a challenge: the untransmitted coefficient of a round
polynomial is $\mathit{claim} + r·(c_1 + … + c_d)$ (\code{transcript.rs:296-302}). Every honest
prover, Flock's included, sends the coefficients of the true round polynomial. The pinned Rust
prover of Flock does not: it derives $G(0) = (\mathit{claim} + r·G(1))·(1 + r)^{-1}$
(\code{zerocheck.rs:116-118}), and at $r = 1$ emits a proof its verifier rejects. That is a
completeness defect of the implementation, of probability at most $(k_{\mathrm{batch}} +
1)/|\E|$; it is recorded in \file{docs/leanvm-target.md} and is no part of \lean{flockError}, a
knowledge-soundness error. Witness: \lean{Phase.Complete} is perfect completeness; the Flock
phase's test runs the round at $r = 1$ with $G(0) ≠ 0$, accepted for the true coefficients and
rejected for the derived ones.}

\reason The inverse is taken by the Rust prover alone; the protocol of the specification, with
a prover that sends the true coefficients, is perfectly complete. Filing its failure under
\lean{flockError}, a knowledge-soundness error, is a category mistake (the dossier's conclusion 3).
\end{change}
The same sentence must reach issue \#3, whose text plans \enquote{exceptional-challenge
completeness conditions} (draft for the tracker only; nothing was posted).
\end{finding}

% ---------------------------------------------------------------- 8.5
\begin{finding}{minor}{f:flock-ring-security-extends-complete}{Knowledge soundness is bundled with
perfect completeness}
\fhead{Evidence} An avoidable coupling; it is what turns
\cref{f:flock-ring-test20-completeness} into a threat to the soundness theorem.
\code{ToArkLib/Component.lean:90-93} (\lean{Security … extends Complete}); the composition uses
only the guard of the first verifier:
\begin{leancode}
  rbr := fun init impl ↦ rbrKnowledgeSoundnessWorstCaseWith_of_eq _
    (Verifier.append_rbrKnowledgeSoundnessWorstCaseWith_of_guarded_first S₁.guarded
      (S₁.kSF init impl) (S₂.kSF init impl) (S₁.rbr init impl) (S₂.rbr init impl))
\end{leancode}
\src{LeanerVM/Protocol/ToArkLib/Component.lean:182-184 at b435631}

\fhead{Classification} A deliberate deviation of the spine from ArkLib's separate statements,
whose reason (one structure per hole) does not need the coupling.

\fhead{Proposed change} A structure \lean{Component.Guarded} with \lean{outputPure} and
\lean{guarded}; \lean{Complete} and \lean{Security} both extend it; \lean{piop_rbrKnowledgeSoundness}
then holds for a protocol whose completeness is unproved or imperfect. (The citation check adds:
the composed \lean{Security} also builds \lean{toComplete := S₁.toComplete.append S₂.toComplete},
\code{ToArkLib/Component.lean:176}, so the proposed split must also give \lean{Security.append} a
completeness-free path; the excerpt above shows only the \lean{rbr} field.)
\begin{change}{convention \emph{Holes}}
\asis (\code{bp:332}) \enquote{\lean{X.Security} (perfect completeness, and round-by-round
knowledge soundness \dots)}
\asproposed \enquote{\lean{X.Security} (the guard, and round-by-round knowledge soundness \dots)}
\reason \lean{piop_rbrKnowledgeSoundness} then holds for a protocol whose completeness is
unproved or imperfect.
\end{change}
\end{finding}

% ---------------------------------------------------------------- 8.6
\begin{finding}{minor}{f:flock-ring-aux-constant-one}{The auxiliary predicate must include the
constant-one position}
\fhead{Evidence} Imprecise wording on a load-bearing definition. \code{c-flock:23-28} \enquote{We
also enforce the position $512$ of every block to be $1$ (preventing the all-zero witness)};
\code{hash.rs:1172-1183}; against \code{bp:850-851} \enquote{\lean{aux q} says that Flock's R1CS
holds of the bits packed into the \code{q_flock} region}, \code{st:239-240},
\code{Spine/Instance.lean:24-25, 143-145}. Read literally, \lean{aux} holds of the zero region and
the lemma \enquote{\lean{aux} $⇒$ the limb slots compress} is false. (The citation check adds that
the blueprint's citation \enquote{Annex C.1} does point at the subsection that contains both
conditions, so the omission is of wording, not of reference.)

\fhead{Classification} An error of the blueprint (of wording).

\fhead{Proposed change}
\begin{change}{the auxiliary predicate, everywhere the phrase occurs}
\asis \enquote{\lean{aux q} says that Flock's R1CS holds of the bits packed into the
\code{q_flock} region}
\asproposed \enquote{\lean{aux q} says that, for every compression, the bits packed into the
\code{q_flock} region satisfy Flock's R1CS \textbf{and hold 1 at the constant position 512} (Annex
C.1, equations (flock:eq:r1cs) and (flock:eq:const-one))}. Add to the acceptance tests:
\enquote{\textbf{The constant position is part of the predicate.} The zero block satisfies the R1CS
and is no compression. Witness: \lean{aux} fails on the zero region.}
\reason Read literally, \lean{aux} holds of the zero region and the lemma \enquote{\lean{aux}
$⇒$ the limb slots compress} is false.
\end{change}
\end{finding}

% ---------------------------------------------------------------- 8.7
\begin{finding}{minor}{f:flock-ring-written-from-spec}{Flock is listed as written from the
specification; its circuit, layout and constants exist only in the Rust and the Python}
\fhead{Evidence} \code{bp:192} \enquote{Flock | A, owned by \#3 | specification Annex C;
\code{crates/flock/src/}}. What the specification does not contain, with the place that fixes it:

\begin{small}
\setlength{\tabcolsep}{4pt}
\begin{tabularx}{\linewidth}{@{}>{\raggedright\arraybackslash}X>{\raggedright\arraybackslash}X@{}}
\toprule
Missing from the specification & Fixed by \\
\midrule
The BLAKE2s circuit and the position of each wire
  & \code{hash.rs:40-64, 139-162, 295-359}; \code{gf2.rs:23-31, 109-154};
  \code{verifier.py:1180-1295} \\ \addlinespace[2pt]
The slot of each limb
  & \code{hash_flock.rs:87-115}; \code{verifier.py:847-850} \\ \addlinespace[2pt]
The selector of a limb claim (a low-index selection; §4.1 has aligned blocks only)
  & \code{stack_open.rs:84-97}; \code{verifier.py:288-302} \\ \addlinespace[2pt]
$g_0$ (\enquote{is public}, \code{c-flock:54})
  & \code{univariate_skip_optimized.rs:104-106}; \code{verifier.py:1095} \\ \addlinespace[2pt]
$φ_8$: which of the eight embeddings, and the modulus of $\F_{2^8}$ (\code{c-flock:35})
  & \code{gf2_8.rs:13}; \code{phi8_tower.rs:15-24}; \code{verifier.py:1092} \\ \addlinespace[2pt]
$k_{\mathrm{batch}} = τ_{\mathrm{BLAKE2S}}$
  & \code{cpu/mod.rs:721, 763}; \code{verifier.py:1405} \\ \addlinespace[2pt]
The floor $τ_{\mathrm{BLAKE2S}} ≥ 3$ (\code{06:80} has the caps only)
  & \code{cpu/mod.rs:162-166}; \code{hash.rs:283-286}; \code{verifier.py:861} \\ \addlinespace[2pt]
The wording \enquote{live in $\qflock$, not the stack} (\code{08:77}): \code{q_flock} is a region
  of the stack
  & \code{cpu/layout.rs:17-25} \\
\bottomrule
\end{tabularx}
\end{small}

\fhead{Classification} A deviation forced by the source (the specification is incomplete); the
blueprint follows the Rust, rightly, and should say so.

\fhead{Proposed change}
\begin{change}{the row \emph{Flock} of the table of artifacts}
\asis \enquote{Flock | A, owned by \#3 | specification Annex C; \code{crates/flock/src/}}
\asproposed \enquote{Flock: protocol and errors | A, owned by \#3 | specification Annex A, Annex C}
and a second row \enquote{Flock: the BLAKE2s circuit, wire positions, limb slots, $φ_8$, the fixed
coordinates, \code{R1CS_DIGEST} | B | \code{crates/flock/src/{hash,gf2}.rs},
\code{zerocheck/univariate_skip_optimized.rs:65, 104-106},
\code{primitives/src/field/phi8_tower.rs:15-24}, \code{lean_vm/src/hash_flock.rs:87-115};
\code{verifier.py:847-850, 1092-1100, 1180-1295}}.
\reason The specification is incomplete; the blueprint follows the Rust, rightly, and should
say so.
\end{change}
Each omission is a line for \file{docs/leanvm-target.md} and for a report to leanVM (drafted here
only).
\end{finding}

% ---------------------------------------------------------------- 8.8
\begin{finding}{minor}{f:flock-ring-test23-ring-switching}{The Flock error quoted in acceptance
test 23 leaves out ring switching}
\fhead{Evidence} \code{bp:1331-1332}: \enquote{$Σ\,\mathit{piopError} < 2^{-150}$ plus Flock's
$< 2^{-183}$}. $2^{-183}$ is $(4·k_{\mathrm{batch}} + 163)/|\E|$ alone (\code{c-flock:282});
\lean{flockError_le} (\code{bp:1085}) adds $2^{32}/|\E| = 2^{-160}$, which dominates by a factor
$2^{23}$.

\fhead{Proposed change}
\begin{change}{acceptance test 23}
\asis \enquote{plus Flock's $< 2^{-183}$}
\asproposed \enquote{plus the Flock phase's, below $2^{-183} + 2^{-160}$, the second term being
ring switching's}
\reason \lean{flockError_le} adds $2^{32}/|\E| = 2^{-160}$, which dominates $2^{-183}$ by a
factor $2^{23}$.
\end{change}
The convention \emph{Errors} wants the closed form next to the theorem: give \lean{flockError} per
challenge as in the table of the dossier's section 5.2.
\end{finding}

% ---------------------------------------------------------------- 8.9
\begin{finding}{minor}{f:flock-ring-ledger-ring-switching}{The ledger row on ring switching
suggests a missing profile; ArkLib's construction is another protocol}
\fhead{Evidence} \code{bp:267}: \enquote{A9 ring switching packing leaves |
\code{RingSwitching/Packing} leaves are \code{sorry}, and no \code{GF(2) → GF(2^64)} profile}. The
dossier's section 6, as its summary puts it: ArkLib's \code{RingSwitching/Packing} is a different protocol
(Diamond--Posen: a carrier message, a batching vector, a relocation sumcheck, error $κ/|L|$),
with completeness and soundness both admitted at the pin and on \code{origin/main}; no
\enquote{profile} turns it into leanVM's. With every leaf proved and a profile for $\F_2 ⊂ \K$,
ArkLib's reduction would still send a carrier element, run a relocation sumcheck, have the point
in $\K$ and the error $κ/|L|$. (The citation check notes that since the merge of \code{144c5aa}
the ArkLib pin is \code{7653a901}, not \code{dca90385}; the 14 \code{sorry} of
\code{RingSwitching/Packing/} are the same at both, so \enquote{admitted at the pin and on
\code{origin/main}} holds at the old and at the new pin.)

\fhead{Classification} An error of the blueprint (of description).

\fhead{Proposed change}
\begin{change}{the ledger row on ring switching}
\asis \enquote{A9 ring switching packing leaves | \code{RingSwitching/Packing} leaves are
\code{sorry}, and no \code{GF(2) → GF(2^64)} profile | Layer 9 | Owned by \#3 (F5); consumed
through the Flock interface.}
\asproposed \enquote{A9 ring switching | ArkLib's \code{Packing} is Diamond--Posen's interactive
reduction (carrier message, batching vector, relocation sumcheck), 14 \code{sorry} at the pin and
on \code{main}; leanVM's is rectangular ($\K ⊗ \E$), batches by a linear map from six challenges
and hands a weighted claim to the commitment scheme | Layer 9 | New work of \#3 (its gate F5); at
most the tensor algebra is shared}.
\reason With every leaf proved and a profile for $\F_2 ⊂ \K$, ArkLib's reduction would still
send a carrier element, run a relocation sumcheck, have the point in $\K$ and the error $κ/|L|$.
\end{change}
\end{finding}

% ---------------------------------------------------------------- 8.10
\begin{finding}{minor}{f:flock-ring-status-python-caps}{The status file's finding that the Python
verifier omits the caps is false at the pin}
\fhead{Evidence} Stale text; it asks for a report upstream that would be wrong.
\code{st:316-319}: \enquote{F9 the Python verifier omits the caps $\mathtt{log\_mem} ∈ [16, 32]$,
$τ_j ≤ 32$, the bytecode power-of-two bound and $τ_{\mathrm{BLAKE2S}} ≥ 3$
(\code{verifier.py:1372-1379} \dots)}. At the pin:
\begin{srccode}
    log_bytecode = log2_strict(len(bytecode)) - BUS_BITS
    require(
        16 <= log_memory <= 32
        and all(0 <= log_height <= 32 for log_height in table_log_heights)
        and table_log_heights[OP_BLAKE2S] >= 3
        and 0 <= log_bytecode <= 32,
        "invalid announced table sizes",
    )
\end{srccode}
\src{python-verifier/verifier.py:857-864}
and \code{log2_strict} requires a power of two (\code{verifier.py:252-254}). The check is present
at the pin; it was assembled in \code{8a6e7750}, \code{348ca455} and \code{5751a5c7}, all
ancestors of \code{a386121f}. (The citation check of the bus dossier corrected the provenance:
the dossier says the lines were introduced by \code{14fbca8f}, which only renamed
\code{BLAKE2S.opcode} to \code{OP_BLAKE2S}; this dossier's citation check had accepted
\code{14fbca8f} from \code{git log -S} on the renamed string. The conclusion stands.)

\fhead{Proposed change} Delete the finding from the status. (It touches this dossier through the
floor on $τ_{\mathrm{BLAKE2S}}$; the other caps belong to another dossier.)
\end{finding}

% ---------------------------------------------------------------- 8.11
\begin{finding}{minor}{f:flock-ring-verifier-needs-flock-def}{The executable verifier needs the
Flock phase's definition, which no roadmap has started}
\fhead{Evidence} Planning. Every proof carries a Flock transcript: the floor is 8 compressions
(\code{hash.rs:283-286}), so the fixture of Layer 12, \enquote{a proof dumped from the pinned Rust
prover \dots\ accepted by \lean{verify}} (\code{bp:1237-1238}), cannot pass before a Lean Flock
verifier exists, with the circuit walk and the constants. The hole for the transcript and
\lean{verify} (K3) lists \enquote{S (schedules, phase \lean{Def}s)} among what it consumes
(\code{bp:614}); the Flock \lean{Def} is \enquote{the inhabitant is \#3's} (\code{bp:609}), and
\#3's first gate has not started.

\fhead{Proposed change} Split the hole of the Flock phase (P6) as every other phase is split. The
definition with its completeness (the schedule of \cref{f:flock-ring-interface-limb-claims}, the
verifier, the circuit walk transcribed, the constants in \file{Parameters}, an honest run on a toy
circuit) is a hole that can be taken now and tested against the Python verifier's walk; the
security stays \#3's. State in the hole for the transcript and \lean{verify} (K3) that its fixture
waits on it.
\end{finding}

% ---------------------------------------------------------------- 8.12
\begin{finding}{note}{f:flock-ring-rust-prover-incomplete}{The pinned Rust prover of Flock is not
perfectly complete}
\fhead{Evidence} A finding against the source, for \file{docs/leanvm-target.md}. The dossier's
section 4.1 and its probe 9.2 (see the evidence of \cref{f:flock-ring-test20-completeness}).

\fhead{Classification} The Lean honest prover \textbf{deliberately deviates} from the Rust prover;
the reason is that the specification's protocol is perfectly complete and the Rust's shortcut is
not; the theorem it owes is \lean{baseProver_complete} for the Lean \lean{prove} (not for the
Rust), and the fixture \enquote{a Rust-produced proof is accepted} is unaffected, the verifiers
being the same.
\end{finding}

% ---------------------------------------------------------------- 8.13
\begin{finding}{note}{f:flock-ring-list-size}{The error for a fixed committed polynomial is not
the error after compilation}
\fhead{Evidence} For the dossier on the compiled verifier. The commitment is list binding only
(\file{b-polynomial-commitment-scheme.tex:4}), so \code{c-flock:264} \enquote{The commitment
determines $\qflock$} holds up to a list. The Rust's parameter analysis multiplies the ring
switching degree by the list size: \code{whir_config.rs:540-541} \enquote{A degree-$d$ identity
test unioned over a Johnson list of size $L$ fails with probability at most $dL/|F|$}, and
\code{:564-567}. Annex A's $2^{-160}$ and the blueprint's \lean{flockError} have no such factor,
which is right in the ideal oracle model; \lean{niError} (\code{bp:1234-1235}) must have it, for
every challenge of every phase. Issue \#3 names the obligation (\enquote{Joint candidate lists}).
\end{finding}

% ---------------------------------------------------------------- 8.14
\begin{finding}{note}{f:flock-ring-seed-constant}{The Fiat--Shamir seed binds a constant that
cannot be recomputed at the pin}
\fhead{Evidence} Extends the status file's finding on the seed (F14). \code{08:55} says the seed
covers \enquote{Flock's BLAKE2s R1CS matrices}. The Rust binds \code{R1CS_DIGEST}, of which
\code{hash.rs:259-275} says: \enquote{baked as an opaque constant \dots\ That recipe needed the
materialized matrices, which this module no longer builds \dots\ To recompute it, check out the
last commit that still had \code{build_matrices}}. The Python copies the bytes
(\code{verifier.py:629}). The constant is transcribed from its source (the blueprint's Category B)
and has no owner in the blueprint; no theorem depends on its value.
\end{finding}

% ---------------------------------------------------------------- 8.15
\begin{finding}{note}{f:flock-ring-stale-comments}{Stale comments in the Rust on what binds the
counter and the flags}
\fhead{Evidence} \code{cpu/mod.rs:568-570} says \enquote{Message, chaining-value, and output words
bind through the memory bus; counter and flags bind through bytecode}, and \code{cpu/mod.rs:619-620}
says \enquote{bytecode binds the counter and flags}. (The citation check corrected the quotation:
the dossier attributed the words \enquote{bytecode binds the counter and flags} to both places;
they are at \code{:619-620} only.) The table reads the metadata cell from memory
(\code{tables.rs:904-907} \enquote{The metadata rides the memory bus like every other operand}),
as the specification says (\code{07:120, 135}). No effect on the protocol; the Lean should cite the
table.

The citation check found two more stale comments of the same kind: \code{zerocheck.rs:389-390}
says \enquote{The prover sends \code{(G(1), G(∞))}}, but the stream carries the coefficients $c_1
= G(0) + G(1) + G(∞)$ and $c_2 = G(∞)$ (\code{zerocheck.rs:118} with \code{transcript.rs:63-71});
\code{hash_flock.rs:8-9} speaks of \enquote{the reduction's two tower-field claims}, while one
ring-switched claim is built (\code{cpu/mod.rs:767}, \code{hash.rs:801-808}).
\end{finding}

% ---------------------------------------------------------------- 8.16
\begin{finding}{note}{f:flock-ring-weight-cube-values}{A weighted claim carries the $2^μ$ cube
values of its weight}
\fhead{Evidence} For the dossier on the opening. \lean{Weight} (\code{Spine/Seams.lean:103-109})
has the field \lean{onCube : CMlPolynomialEval E μ}, a vector of $2^μ$ elements of $\E$. The Flock
verifier's output statement contains one. As a specification this is fine; as a function that
runs it builds a table of $2^μ$ values ($μ$ up to 28) that no deployed verifier builds
(\code{ring_switch.rs:63-66}). The field \lean{mle_eq}, for the ring switching weight, is the boxed
identity of \code{a-ring:150-154}, owed by the Flock roadmap.
\end{finding}
